
\section{\BookRecordIndex{T0408}正则性判据}\label{sec:3-3}
\zhindex{zhengzexingpanju}{正则性判据}

在继续发展一般理论之前，我们先举几个例子说明\BookRecordIndex{T0408}正则性的概念。在可积情形以及
$\Sigma=S^2=\mathbb{CP}^1$ 时，这一概念尤其清楚。Grothendieck \srcite{162} 证明，
$\mathbb{CP}^1$ 上任意全纯向量丛 $E$ 都全纯等价于若干全纯线丛的直和，而且这种分裂
在不计各直和项次序的意义下唯一。因此
\[
  E=L_1\oplus\cdots\oplus L_n
\]
完全由陈数集合 $c_1(L_1),\ldots,c_1(L_n)$ 刻画。这里及下文中，我们把陈类 $c_1(L)$
与相应陈数 $\langle c_1(L),[S^2]\rangle$ 等同。特别地，当 $J$ 是可积复结构且
$u:S^2\to M$ 是 $J$-全纯球面时，Grothendieck 定理可用于向量丛 $E=u^*TM$。注意，
总和
\[
  c_1(E)=\sum_i c_1(L_i)
\]
是拓扑不变量；但当 $u:S^2\to M$ 在 $J$-全纯球面空间的某个连通分支中变化时，分解
$E=u^*TM$ 中的各个数 $c_1(L_1),\ldots,c_1(L_n)$ 可能随之变化。

下面的结果立即来自定理 C.1.10\textup{(iii)}；该定理断言，
$\mathbb{CP}^1$ 上复线丛的实线性\BookRecordIndex{T0210}柯西--黎曼算子或者单射，或者满射，或者二者兼具。
尽管附录 C 给出了完整证明，我们仍在这里简述一个基于代数几何观察的论证，以说明建立在
椭圆型偏微分方程上的分析框架，如何与全纯向量丛的经典结果相契合。

\begin{lemma}\label{lem:3-3-1}
设 $J$ 可积，$u:\mathbb{CP}^1\to M$ 是 $J$-全纯球面。若 $u^*TM$ 的每个直和项的
陈数都满足 $c_1\ge-1$，则 $D_u$ 是满射。
\end{lemma}

\begin{proof}
由于 $J$ 可积，算子 $D_u$ 复线性，并且恰是由 $M$ 的全纯坐标卡确定的 Dolbeault 微分
$\bar\partial$。因此，它保持 $u^*TM$ 的全纯线丛直和分裂，可以分别考察
$u^*TM$ 的每个全纯直和项\sourcefootnote{原书此处写作“$(TM,J)$ 的每个线子丛 $L$”；然而该丛的底空间是
$M$，而这里的算子及其分裂定义在 $\mathbb{CP}^1$ 上。由紧邻的分裂式和下文
Dolbeault 算子的定义域可知，应为 $u^*TM$ 的全纯直和项。}，记为 $L$。此外，
\[
  D_u=\bar\partial:
  \Omega^0(\mathbb{CP}^1,L)
  \longrightarrow\Omega^{0,1}(\mathbb{CP}^1,L)
\]
的余核恰是 Dolbeault 上同调群
$H_{\bar\partial}^{0,1}(\mathbb{CP}^1,L)$。所以，只须证明当 $c_1(L)\ge-1$ 时这个群
为零。对任意全纯线丛 $L\to\mathbb{CP}^1$，有
\[
  H_{\bar\partial}^{0,1}(\mathbb{CP}^1,L)
  \cong
  \Bigl(H_{\bar\partial}^{1,0}(\mathbb{CP}^1,L^*)\Bigr)^*.
\]
这里，$H_{\bar\partial}^{1,0}(\mathbb{CP}^1,L^*)$ 是取值于对偶丛 $L^*$ 的全纯
$1$-形式空间，因而同构于丛 $L^*\otimes K$ 的全纯截面空间
$H^0(\mathbb{CP}^1,L^*\otimes K)$，其中
$K:=T^*\mathbb{CP}^1$ 是典范线丛。这是 Kodaira--\BookRecordIndex{T0048}Serre 对偶的一个简单特例
（比较 \srcite{159}，第 1 章第 2 节），也可直接考察转移映射来验证。（注意，
$L^*\otimes K$ 在 $z$ 处的纤维由同态 $T_zS^2\to L_z^*$ 组成，也就是 $z$ 处取值于
$L_z^*$ 的 $1$-形式。）

现在，$\mathbb{CP}^1$ 上线丛 $L'$ 有非零全纯截面，当且仅当 $c_1(L')\ge0$。这是
Kodaira 消没定理的一个特例，也可由\BookRecordIndex{T0186}交点正性以及把 $c_1(L')$ 解释为零截面的自交数得到。
因此，$D_u$ 是满射，当且仅当 $E$ 的每个全纯直和项 $L$ 都满足
\[
  c_1(L^*\otimes K)<0.
\]
由于 $c_1(L^*\otimes K)=-c_1(L)-2$，这等价于 $c_1(L)\ge-1$。引理得证。
\end{proof}

上述论证可如下推广到殆复流形。

\begin{lemma}\label{lem:3-3-2}
设 $E\to S^2$ 是秩为 $n$ 的复向量丛，并且
\[
  D:\Omega^0(S^2,E)\longrightarrow\Omega^{0,1}(S^2,E)
\]
是实线性\BookRecordIndex{T0210}柯西--黎曼算子。假设存在\BookRecordIndex{T0239}滤过
\[
  E_1\subset E_2\subset\cdots\subset E_n=E,
\]
其中每个 $E_k$ 都是秩为 $k$ 的 $D$-不变复子丛。令 $E_0:=0$。
则 $D$ 是满射，当且仅当\sourcefootnote{原书此处及证明末句均断言“当且仅当”。
按 $D$-不变子丛的通常含义，充分性成立，但必要性一般不成立：例如
$\mathbb{CP}^1$ 上有非平凡正合列
$0\to\mathcal O(-2)\to\mathcal O(-1)\oplus\mathcal O(-1)
\to\mathcal O\to0$；中项的 Dolbeault 算子满射，而\BookRecordIndex{T0239}滤过的第一个商的陈数为 $-2$。
这里保留原书断言；下文引理~\ref{lem:3-3-3} 的\BookRecordIndex{T0239}滤过中，第一个子丛上的限制算子本身
满射，此时由相应的核余核正合列可知，总算子满射当且仅当商算子满射，故该应用不受此
反例影响。}对每个 $k$ 都有
\[
  c_1(E_k/E_{k-1})\ge-1.
\]
\end{lemma}

\begin{proof}
对 $n$ 作归纳。当 $n=1$ 时，断言来自定理 C.1.10\textup{(iii)}。设 $n\ge2$，并假定
结论对 $n-1$ 成立。算子 $D$ 在 $E_{n-1}$ 的截面上的限制
\[
  D':\Omega^0(S^2,E_{n-1})
  \longrightarrow\Omega^{0,1}(S^2,E_{n-1})
\]
是满射，当且仅当对 $k=1,\ldots,n-1$ 都有
$c_1(E_k/E_{k-1})\ge-1$。另一方面，由定理 C.1.10\textup{(iii)}，$D$ 确定的商算子
\[
  D'':\Omega^0(S^2,E_n/E_{n-1})
  \longrightarrow\Omega^{0,1}(S^2,E_n/E_{n-1})
\]
是满射，当且仅当 $c_1(E_n/E_{n-1})\ge-1$。因此，归纳步骤来自如下事实：$D$ 是满射，
当且仅当 $D'$ 和 $D''$ 都是满射。引理得证。
\end{proof}

\begin{lemma}\label{lem:3-3-3}
设 $J$ 是四维流形 $M$ 上的\BookRecordIndex{T0088}殆复结构，
$u:\mathbb{CP}^1\to M$ 是浸入的 $J$-全纯球面。则 $D_u$ 是满射，当且仅当
\[
  c_1(u^*TM)\ge1.
\]
\end{lemma}

\begin{proof}
由式~\eqref{eq:3-1-4} 中 $D_u$ 的定义，对每个向量场 $\zeta\in\Vect(\Sigma)$ 以及
每条 \BookRecordIndex{T0003}$J$-全纯曲线 $u$，都有
\[
  D_u(du\circ\zeta)=du\circ\bar\partial_j\zeta.
\]
若 $u$ 是浸入，就可知复子丛
\[
  L:=\operatorname{im}du\subset u^*TM
\]
在 $D_u$ 下不变。当 $\Sigma=\mathbb{CP}^1$ 时，$c_1(L)=2$，并且
\[
  c_1(u^*TM/L)=c_1(u^*TM)-2.
\]
这两个陈数都不小于 $-1$，当且仅当 $c_1(u^*TM)\ge1$。所以结论来自
引理~\ref{lem:3-3-2}。
\end{proof}

\begin{corollary}\label{cor:3-3-4}
设 $J$ 是四维流形 $M$ 上的\BookRecordIndex{T0088}殆复结构，$C$ 是自交数为 $C\cdot C=p$ 的嵌入
$J$-全纯球面。则 $J$ 对同调类 $A=[C]$ 正则，当且仅当 $p\ge-1$。
\end{corollary}

\begin{proof}
由定理~\ref{thm:2-6-4}，\BookRecordIndex{T0246}模空间 $\mathcal M(A;J)$ 中的曲线全都嵌入。因此结论来自
引理~\ref{lem:3-3-3}。
\end{proof}

\begin{corollary}\label{cor:3-3-5}
令 $\widetilde M:=S^2\times M$ 是 $S^2$ 与\BookRecordIndex{T0358}辛流形 $(M,\omega)$ 的乘积，并令
$\widetilde A\in H_2(\widetilde M;\mathbb Z)$ 是由球面 $S^2\times\{\mathrm{pt}\}$
表示的同调类。则对每个 $J\in\mathcal J(M,\omega)$，乘积\BookRecordIndex{T0088}殆复结构
$\widetilde J:=i\times J$ 都对 $\widetilde A$ 正则。
\end{corollary}

\begin{proof}
每个 $\widetilde J$-全纯 $\widetilde A$-球面都形如
\[
  \widetilde u(z)=(\phi(z),x_0),
\]
其中 $\phi$ 是\BookRecordIndex{T0126}分式线性变换，$x_0\in M$。因此 $D_{\widetilde u}$ 复线性，而且拉回
切丛\sourcefootnote{原书此处写作
$\widetilde u^*TM$；但 $\widetilde u$ 的靶是 $\widetilde M=S^2\times M$，且紧随其后的
分解包含来自 $S^2$ 因子的陈数为 $2$ 的切向线丛，故改为
$\widetilde u^*T\widetilde M$。}$\widetilde u^*T\widetilde M$ 分解成
$D_{\widetilde u}$-不变线丛的直和；其中一个线丛的
陈数为 $2$，其余线丛的陈数均为零。由引理~\ref{lem:3-3-2}，
$D_{\widetilde u}$ 是满射。
\end{proof}

\begin{example}\label{ex:3-3-6}
考虑乘积流形
\[
  M=S^2\times S^2
\]
及其乘积\BookRecordIndex{T0088}殆复结构 $J_0=j\times j$。到两个因子的投影 $\pi_i$ 都是全纯映射，所以任意
全纯映射 $u:S^2\to S^2\times S^2$ 都形如
\[
  z\longmapsto\bigl(v_1(z),v_2(z)\bigr),
\]
其中 $v_i:S^2\to S^2$ 全纯且次数 $d_i$ 非负。于是 $u^*TM$ 全纯分裂为
$L_1\oplus L_2$，其中 $L_i:=v_i^*TS^2$ 的次数为 $2d_i$。所以这些曲线全都正则。
然而，$S^2\times S^2$ 还带有其他被\BookRecordIndex{T0363}辛形式\BookRecordIndex{T0372}驯服的复结构，并且其中存在非正则曲线。

看清这一点的最佳方式，是考虑如下得到的 Hirzebruch 复结构 $J_{2k}$：把
$S^2\times S^2$ 等同于 $\mathbb P(L\oplus\mathbb C)$，其中 $L$ 是次数
$2k>0$ 的全纯线丛。子丛 $L$ 对应于 $\mathbb P(L\oplus\mathbb C)$ 的一个截面
$C_L$，其\BookRecordIndex{T0117}法丛为 $L^*$。所以当 $k>0$ 时，$c_1(L^*)<-1$，从而 $C_L$ 不正则。可以
选取 $\mathbb P(L\oplus\mathbb C)$ 与 $S^2\times S^2$ 的一个等同，把 $C_L$ 映到
反对角线
\[
  \bigl(z,A(z)\bigr),
\]
其中 $A$ 是对径映射。（考察局部平凡化即可显式作出这一等同。）此外，
$\mathbb P(L\oplus\mathbb C)$ 上存在 凯勒 形式，在该等同下被映为\BookRecordIndex{T0363}辛形式
\[
  \omega^\lambda:=\lambda\pi_1^*(\sigma)\oplus\pi_2^*(\sigma),
\]
其中 $\sigma$ 是 $S^2$ 上的面积形式，$\lambda>1$。总之，当 $\lambda>1$ 时，\BookRecordIndex{T0358}辛流形
$(S^2\times S^2,\omega^\lambda)$ 带有某个与 $\omega^\lambda$ \BookRecordIndex{T0348}相容的 $J$，使反对角线
成为一条非正则曲线。（例如可见 \srcite{5}，其中把这些流形构造成约化空间；
\srcite{262} 对作为\BookRecordIndex{T0358}辛流形的直纹曲面作了广泛讨论。）

构造非正则曲线的另一种方法是吹起。从 $\mathbb{CP}^2$ 出发作两次吹起，并让第二个吹起点
落在第一次吹起所得 $X_1$ 的例外除子 $C$ 上。由于 $C$ 的\BookRecordIndex{T0117}法丛陈类为 $-1$，它在单点
吹起 $X_1$ 中正则。然而，它在双点吹起 $X_2$ 中的严格变换 $C'$ 的\BookRecordIndex{T0117}法丛陈类为 $-2$，
所以并不正则。事实上，也可在反对角线之外的一点吹起 Hirzebruch 曲面
$(S^2\times S^2,J_2)$，得到这个例子。不过，从后一幅图景很容易看出，非正则曲线怎样
在 $J$ 变化时消失（比较注~\ref{rem:3-2-8}）：只要把第二个吹起点从例外除子 $C$ 上
移开，$X_2$ 中就不再有\BookRecordIndex{T0117}法丛陈类为 $-2$ 的曲线。
\end{example}

这是我们的讨论中第一次必须把曲线限制为球面。其他结果适用于定义域 $(\Sigma,j)$ 固定的
\BookRecordIndex{T0003}$J$-全纯曲线，但这类曲线几乎从不正则，因为即使 $J$ 可积，条件
\[
  H_{\bar\partial}^{0,1}(\Sigma,u^*TM)=0
\]
通常也不成立。事实上，若 $u$ 嵌入，且其像 $C:=u(\Sigma)$ 的切丛 $TC$ 是
$\left.TM\right|_C$ 的一个全纯直和项，则 $H_{\bar\partial}^{0,1}(\Sigma,u^*TM)$
包含非平凡空间
\[
  H_{\bar\partial}^{0,1}(\Sigma,T\Sigma),
\]
即 $\Sigma$ 的 Teichmüller 空间在 $j$ 处的切空间。因此，当 $J$ 在 $M$ 上变化时，为了
找到曲线的相应形变，通常也必须让 $\Sigma$ 上的 $j$ 变化。为此，需要建立允许 $j$ 在
Teichmüller 空间中变化的弗雷德霍姆理论。这在当前阶段不会造成困难，却会使紧致性的讨论
更复杂，因为 Teichmüller 空间即使对映射类群取商也不紧致。

上面讨论中的许多部分仍适用于高亏格嵌入曲线。例如，如果已知全纯向量丛 $u^*TM$ 是线丛
之和，引理~\ref{lem:3-3-2} 中的论证仍然适用。不过需要稍加小心：对于亏格为 $g$ 的
曲线上的线丛，是否存在全纯截面不再仅由陈数这一拓扑量决定，除非这个数大于 $2g-2$
（见定理 C.1.10）。在四维嵌入情形，对线丛 $u^*TM/TC$ 的这一条件等价于要求
$c_1(A)>0$。

在四维中，代表类 $A$ 的嵌入 \BookRecordIndex{T0003}$J$-全纯曲线的亏格由\BookRecordIndex{T0055}伴随公式
\[
  2-2g+A\cdot A=c_1(A)
\]
确定。对任意亏格，只要 $c_1(A)>0$，就能证明这个\BookRecordIndex{T0246}模空间正则。该分析论证最早由
\personlater{geluomofu}{格罗莫夫}{Gromov}提出，Hofer--Lizan--Sikorav \srcite{179}
完成。先用\BookRecordIndex{T0186}交点正性证明\BookRecordIndex{T0246}模空间完全由
嵌入曲线组成，再像引理~\ref{lem:3-3-3} 的证明那样论证。这里如上所述，必须重述理论，
让 $\Sigma$ 上的复结构 $j$ 在 Teichmüller 空间中自由变化。这样，论证归结为证明：黎曼
曲面上负次数复线丛的实线性\BookRecordIndex{T0210}柯西--黎曼算子总是单射（见定理 C.1.10\textup{(iii)}）。
对浸入球面，也可用几何比较论证得到类似结果，见 McDuff \srcite{260}。
Ivashkovich--Shevchishin \srcite{202} 把 \srcite{179} 的讨论推广到了颇有意思的奇异曲线
情形。

下一个例子说明考察高亏格嵌入曲线时出现的一些新特征。

\begin{example}[\BookRecordIndex{T0328}椭圆曲面]\label{ex:3-3-7}
考虑 $n\ge1$ 时的\BookRecordIndex{T0328}椭圆曲面 $E(n)$。$E(1)$ 就是 $\mathbb{CP}^2$ 的九点吹起；选取这九点，
使 $E(1)$ 带有一个到射影直线的 \BookRecordIndex{T0033}Lefschetz 纤维化，其一般纤维是椭圆曲线。为此，只须把
这九点选为 $\mathbb{CP}^2$ 中两条不同三次曲线 $f_1=0$ 与 $f_2=0$ 的交点，因为这些点
都落在曲线族
\[
  \Gamma_\lambda
  :=\{\lambda_1f_1+\lambda_2f_2=0\},
  \qquad
  [\lambda_1:\lambda_2]\in\mathbb{CP}^1
\]
上（例如见 \srcite{277}，例 7.6）。把这九点吹起后，各曲线 $\Gamma_\lambda$ 彼此分开，
因而构成一个到 $S^2$ 的映射的纤维；一般而言，这是以环面为纤维的光滑纤维丛，但在使
$\Gamma_\lambda$ 奇异的有限多个 $\lambda$ 上有特殊纤维。一般地，$E(n)$ 定义为
$E(1)$ 的 $n$ 重分歧覆盖；在辛语境中，这与 $n$ 重\BookRecordIndex{T0278}纤维和
\[
  E(1)\mathbin{\#_F}\cdots\mathbin{\#_F}E(1)
\]
相同。特别地，$E(2)$ 是 K3 曲面，见 \srcite{277}，习题 7.31。

考虑纤维类 $F$，以及计数代表这个类的嵌入环面这一问题。若建立模问题时允许定义域上的
复结构 $j_\Sigma$ 变化，则其指标为零。具体地，相关\BookRecordIndex{T0210}柯西--黎曼算子的指标为
\[
  2n(1-g)+2c_1(A)+2,
\]
其中直和项 $2$ 是 Teichmüller 空间的维数；再减去环面的全纯自同构群（平移群）的维数
$2$。此外，可以改写引理~\ref{lem:3-3-1} 的证明，得到：复曲面中的嵌入环面正则，
当且仅当其\BookRecordIndex{T0117}法丛
\[
  L=u^*TM/T\Sigma
\]
的对偶没有截面。当 $c_1(L)=0$ 时，这等价于 $L$ 非平凡。

Kodaira 对复曲面的分类表明，当 $n>2$ 时，$E(n)$ 上每个复结构都容许一个到 $S^2$ 的
\BookRecordIndex{T0279}纤维化，其一般纤维是环面。（细节见 Beauville \srcite{33}，第 IX 章。）所以，对可积
$J$，$F$ 类中总有一个实二参数的嵌入环面族，它们并不正则。相反，尽管
$T^2\times S^2$ 上的乘积结构对 $F=[T^2\times\mathrm{pt}]$ 不正则，这个流形确实带有
对 $F$ 正则的复结构：可以把它等同于直纹曲面 $\mathbb P(L\oplus\mathbb C)$，其中
$L\to T^2$ 是满足 $c_1(L)=0$ 的非平凡丛。（这个直纹曲面恰有两个 $F$ 类正则全纯
环面，即 $\mathbb P(L\oplus\{0\})$ 和 $\mathbb P(\{0\}\oplus\mathbb C)$。）

这些情形的一个重要区别是：如 McDuff \srcite{265} 所述，$E(n)$ 中 $F$ 类环面的带符号
计数为 $2-n$，所以在 $n>2$ 时为负；$T^2\times S^2$ 中相应的数为正，等于 $2$。由
注~\ref{rem:3-2-6}，$E(n)$（$n>2$）上不存在对 $F$ 正则的复结构 $J$。另一方面，$J$
的一个通有（不可积）扰动容许有限多个嵌入环面，它们的符号之和为 $2-n$。这也可用\BookRecordIndex{T0399}障碍丛
方法解释。
\end{example}
